\chapter{Techno-Economic Analysis}
\label{ch:techno_economic}

To evaluate the commercial viability and financial feasibility of the custom electromagnetic brake joint for a Portable Coordinate Measurement Arm (PCMA), a techno-economic analysis was performed. The evaluation assesses actual prototype procurement expenditures, university workshop fabrication overheads, scaled industrial manufacturing projections, and lifecycle economics relative to commercial metrology joint locks.

\section{Actual Prototype Expenditure Breakdown}
Table~\ref{tab:actual_project_expenses} itemizes the verified procurement and fabrication expenses incurred during prototype construction, categorized across certified vendor quotes and internal workshop charges[cite: 30, 32, 34, 36, 38].

\begin{table}[H]
	\centering
	\caption{Verified Prototype Component Procurement and Workshop Expenditure}
	\label{tab:actual_project_expenses}
	\small
	\begin{tabular}{lllr}
		\toprule
		\textbf{Vendor / Source} & \textbf{Document Reference} & \textbf{Expense Description} & \textbf{Total (VAT Incl.)} \\
		\midrule
		Fabrinox (Pty) Ltd    & PO US10000352153 & Laser-cut cover, pressure platform \& friction plate & R 214.89 \\
		RS Components South Africa & Basket Quote & Enamelled wire, MOSFETs, diodes, bearings \& couplings & R 3 550.16 \\
		Micro Robotics        & Quote JOB-16439 & 1000PPR optical encoder \& step-down buck regulator & R 802.70 \\
		Takealot              & Retail Receipt & High-temperature Kapton insulation tape & R 299.99 \\
		SU Mechanical Workshop& Invoice (Miss S. Paulse) & Raw material stock (R 3 084.81) + Labor 18.5 hrs (R 12 025.00) & R 15 109.81 \\
		\midrule
		\textbf{Total Project Development Expenditure} & & & \textbf{R 19 977.55} \\
		\bottomrule
	\end{tabular}
\end{table}

\section{Scaled Production Cost Estimation}
While initial prototype development required substantial custom machining labor (18.5 workshop hours at R 650.00/hr) and low-volume retail part ordering[cite: 30], industrial scaling significantly reduces unit costs through automated coil winding, surface-mount PCB assembly (PCBA), and bulk material purchasing. Applying an $85\%$ experience curve model for precision mechatronics across a batch size of 1 000 units yields the scaled unit production estimates summarized in Table~\ref{tab:scaled_costs}.

\begin{table}[H]
	\centering
	\caption{Cost Reduction Comparison Across Production Scales}
	\label{tab:scaled_costs}
	\small
	\begin{tabular}{lrrr}
		\toprule
		\textbf{Cost Category} & \textbf{Prototype (1 Unit)} & \textbf{Medium Batch (100 Units)} & \textbf{Mass Production (1 000 Units)} \\
		\midrule
		Raw Materials \& Metalwork  & R 3 299.70 & R 850.00 & R 480.00 \\
		Electronics, Sensors \& Wire & R 4 652.85 & R 1 120.00 & R 540.00 \\
		Machining \& Assembly Labor  & R 12 025.00 & R 650.00 & R 120.00 \\
		\midrule
		\textbf{Unit Cost Total}     & \textbf{R 19 977.55} & \textbf{R 2 620.00} & \textbf{R 1 140.00} \\
		\bottomrule
	\end{tabular}
\end{table}

\section{Lifecycle Economics and Capital Expenditure Impact}
Commercial PCMA joint assemblies utilizing imported pneumatic or precision electromagnetic brakes typically cost upwards of R 12 000.00 per joint module. At a mass-production cost of approximately R 1 140.00 per joint, the proposed design achieves an $80\text{--}90\%$ reduction in initial hardware capital expenditure (CapEx). Furthermore, by eliminating mechanical wear linkages, pneumatic supply lines, and fluid seals, operational maintenance costs (OpEx) are minimized over an estimated 10-year operating life.